\chapter[Introdução]{Introdução}
%\addcontentsline{toc}{chapter}{Introdução}

\section{Contexto}
Nos últimos anos, especialmente a partir de 2022, os Modelos de Linguagem de Grande Escala, do inglês \textit{Large Language Models} (LLMs), passaram a integrar de forma cada vez mais significativa o processo de desenvolvimento de software. Essas ferramentas têm sido utilizadas em diversas atividades, como geração de código, explicação de trechos, sugestão de correções e apoio à análise de programas. Com isso, os LLMs passaram a ocupar um papel relevante no cotidiano de desenvolvedores e equipes de engenharia de software, ainda que continuem apresentando limitações em tarefas que demandam maior precisão e compreensão da lógica do código \cite{zhou_large_2024}.

Com esse crescimento, também emergiram novas formas de desenvolvimento de software apoiadas por IA. Uma delas é o chamado \textit{vibe coding}, no qual o desenvolvedor descreve, em linguagem natural, o que deseja construir e realiza ajustes sobre o código gerado pelo modelo. Paralelamente, ganhou destaque uma abordagem mais controlada, conhecida como \textit{AI-assisted engineering}, na qual a IA atua como suporte dentro de um processo supervisionado, enquanto o desenvolvedor permanece responsável por revisar, validar e decidir o que deve ser incorporado ao sistema \cite{osmani_vibe_2025}. Essa distinção é relevante, pois evidencia que o impacto da IA no desenvolvimento de software não depende apenas das capacidades do modelo, mas também da forma como ele é integrado ao processo de desenvolvimento.

Apesar dos ganhos de produtividade proporcionados por essas ferramentas, o uso de LLMs em tarefas relacionadas à segurança de software tem gerado preocupações. Estudos indicam que assistentes de programação baseados em IA podem produzir código inseguro em uma parcela significativa dos casos. Em uma pesquisa sobre o GitHub Copilot, por exemplo, cerca de 40\% dos programas gerados nos cenários analisados apresentavam vulnerabilidades, evidenciando que essas ferramentas podem acelerar o desenvolvimento, mas não devem ser utilizadas sem validação, especialmente em contextos críticos \cite{pearce_asleep_2021}.

Outro fator relevante é que o desempenho desses modelos depende fortemente do contexto fornecido. Em tarefas de detecção de vulnerabilidades, pesquisas mostram que a inclusão de informações adicionais pode melhorar a capacidade da IA de diferenciar falhas reais de situações apenas aparentemente suspeitas \cite{iranmanesh_zerofalse_2025}.

Diante disso, trabalhos recentes passaram a explorar abordagens em que a LLM não atua de forma isolada, mas em conjunto com ferramentas externas de apoio. Entre essas ferramentas, destacam-se as soluções de \textit{Static Application Security Testing} (SAST), que analisam o código-fonte sem executá-lo e apontam possíveis vulnerabilidades com base em padrões e regras conhecidas. Embora sejam úteis por ampliarem a cobertura da análise, essas ferramentas indicam condições potencialmente vulneráveis que ainda precisam ser interpretadas e validadas. Assim, seus alertas não devem ser tratados isoladamente como correções ou comprovações definitivas de vulnerabilidade, mas como evidências técnicas que podem orientar etapas posteriores de análise e remediação \cite{iranmanesh_zerofalse_2025}.

Nesse contexto, a integração entre LLMs e ferramentas SAST apresenta-se como uma alternativa promissora para tornar a correção de vulnerabilidades mais eficiente e adequada ao contexto do código analisado. Enquanto as ferramentas SAST contribuem para a identificação de possíveis problemas no código-fonte, os LLMs podem apoiar uma análise mais contextualizada, auxiliando na compreensão e na correção das vulnerabilidades detectadas. Estudos recentes indicam que abordagens iterativas, baseadas no uso de contexto adicional e no feedback de ferramentas externas, tendem a produzir resultados mais confiáveis do que interações isoladas fundamentadas em um único prompt \cite{sriram_improving_2026}. É a partir dessa possibilidade de combinação entre essas tecnologias que se estabelece a motivação deste trabalho.

\section{Problema}

Embora a integração entre ferramentas de segurança e inteligência artificial se mostre promissora, sua aplicação prática ainda apresenta diversos desafios. Um dos principais está relacionado ao fato de que os modelos de linguagem, quando utilizados de forma isolada para corrigir vulnerabilidades, ainda não produzem resultados suficientemente confiáveis. Estudos recentes indicam que esses modelos conseguem identificar a existência de falhas com desempenho razoável, mas ainda apresentam grande dificuldade para gerar correções efetivas. Em um desses trabalhos, por exemplo, o modelo alcançou 76\% de acerto na detecção de vulnerabilidades, mas apenas 15\% de sucesso no reparo quando não recebeu contexto adicional sobre o problema \cite{khan_code_2025}.

Além da baixa taxa de sucesso no reparo, outro problema relevante é que, em muitos casos, a tentativa de corrigir a vulnerabilidade pode comprometer o funcionamento original do software. Em vez de produzir uma solução segura e funcional, o modelo pode remover partes importantes da lógica da aplicação ou gerar um código que, embora elimine a falha de segurança, prejudique a utilidade do sistema. Esse comportamento evidencia que corrigir uma vulnerabilidade não consiste apenas em remover o erro, mas também em preservar o comportamento esperado do programa \cite{dai_rethinking_2025}.

Por outro lado, depender exclusivamente de ferramentas tradicionais de análise estática também não resolve o problema. Embora essas ferramentas sejam importantes para localizar possíveis falhas no código, seus alertas ainda precisam ser interpretados, priorizados e validados por desenvolvedores. Além disso, a identificação de uma \textit{issue} não implica, por si só, a existência de uma correção segura e funcional. Dessa forma, o SAST contribui para apontar evidências técnicas sobre possíveis vulnerabilidades, mas não substitui etapas posteriores de análise, remediação e validação \cite{iranmanesh_zerofalse_2025}.


Pesquisas recentes indicam que o desempenho das LLMs em tarefas de segurança pode melhorar quando o modelo recebe mais contexto e quando seu processo de resposta é conduzido de forma mais estruturada \cite{tao_adapting_2025}. Nesse cenário, estratégias como \textit{Retrieval-Augmented Generation} (RAG), que permite recuperar informações relevantes antes da geração da resposta, e \textit{Chain-of-Thought} (CoT), que incentiva um raciocínio em etapas, surgem como alternativas promissoras para apoiar tanto a interpretação dos alertas quanto o reparo das vulnerabilidades. Ainda assim, não está claro até que ponto essas estratégias realmente aumentam a eficácia da LLM quando ela atua em conjunto com ferramentas SAST \cite{sriram_improving_2026}.

Assim, o problema central desta pesquisa surge justamente das limitações dessas abordagens quando utilizadas de forma isolada. De um lado, as LLMs ainda apresentam dificuldade para reparar vulnerabilidades de maneira confiável e sem comprometer a funcionalidade do software. De outro, ferramentas SAST são capazes de apontar possíveis falhas no código, mas não produzem, por si só, uma correção segura e funcional. Essa situação evidencia uma lacuna importante na engenharia de software atual: compreender até que ponto essas tecnologias podem atuar de forma complementar, com o SAST auxiliando na localização e descrição dos problemas, e a LLM contribuindo para gerar correções que possam ser posteriormente validadas por reanálise estática e testes automatizados.


Com base nisso, esta pesquisa busca responder às seguintes questões:
\begin{itemize}
 \item Q1: É possível gerar correções válidas para vulnerabilidades por meio de um \textit{pipeline} automatizado que integra LLMs, ferramentas SAST e testes automatizados, considerando a remoção dos alertas, a preservação dos testes automatizados e a não introdução de novas \textit{issues} de segurança?
 \item Q2: Como diferentes estratégias de prompt, incluindo prompt simples, raciocínio guiado e enriquecimento de contexto por RAG simplificado, influenciam a taxa de correções válidas geradas pelo modelo de linguagem?
\item Q3: O desempenho do \textit{pipeline} varia de acordo com o tipo de vulnerabilidade ou categoria associada ao OWASP Top 10?
 \item Q4: Quais limitações práticas surgem ao aplicar um \textit{pipeline} automatizado de correção de vulnerabilidades em uma base de código real, com dependências, arquitetura própria e suíte de testes existente?
 \item Q5: Como diferentes modelos de linguagem se comparam na geração de correções válidas pelo \textit{pipeline}, considerando eficácia, custo e tempo de execução?
\end{itemize}

\section{Objetivos}
\subsection{Objetivo Geral}

Investigar a eficácia de um fluxo iterativo que integra Modelos de Linguagem de Grande Escala (LLMs), ferramentas de Análise Estática de Segurança (SAST) e testes automatizados na correção de vulnerabilidades em código-fonte, considerando a remoção dos alertas de segurança, a preservação da funcionalidade original e a não introdução de novas \textit{issues}.


\subsection{Objetivos Específicos}
\begin{itemize}
\item Desenvolver um \textit{pipeline} automatizado que integre uma ferramenta de análise estática de segurança, um modelo de linguagem e testes automatizados para apoiar a correção de \textit{issues} de segurança em código-fonte.

\item Definir e aplicar diferentes estratégias de \textit{prompt}, incluindo \textit{prompt} simples, raciocínio guiado e enriquecimento de contexto por RAG simplificado, a fim de comparar seu impacto na geração de correções válidas.

\item Avaliar as correções geradas a partir da reanálise SAST, da execução dos testes automatizados e da verificação de que novas \textit{issues} de segurança não foram introduzidas.

\item Analisar os resultados obtidos no projeto, considerando a taxa de remediação validada, os tipos de \textit{issues} detectadas, as falhas de correção e as limitações práticas do \textit{pipeline}.

\item Comparar diferentes modelos de linguagem na geração de correções válidas pelo \textit{pipeline}, considerando a eficácia, o custo e o tempo de execução nas condições avaliadas.

\end{itemize}


\section{Metodologia}
Este trabalho realizou uma investigação empírica sobre a eficácia de um \textit{pipeline} que integra Modelos de Linguagem de Grande Escala (LLMs), Análise Estática de Segurança (SAST) e testes automatizados para apoiar a correção de alertas de segurança em código-fonte. A avaliação examinou diferentes estratégias de \textit{prompt} e modelos de linguagem, considerando a remoção dos alertas, a preservação funcional e a ausência de novos alertas nas condições definidas para o estudo.

O experimento foi conduzido no OWASP Juice Shop, uma aplicação web de código aberto e deliberadamente vulnerável. A análise estática com Semgrep forneceu os alertas utilizados para constituir um catálogo de 78 casos. O projeto foi executado em ambiente preparado com Docker, mantendo uma revisão fixada do código e as condições de avaliação registradas nas campanhas.

Foram avaliados os modelos GPT-4.1 mini, GPT-5.6 Luna e GPT-6 Luna sobre os mesmos casos. Para cada caso, as estratégias seguiram a ordem P1, P2, P3 e P4: instrução direta, raciocínio guiado, diagnóstico SAST com documentação recuperada por RAG simplificado e combinação do diagnóstico e da documentação com o raciocínio guiado. Como as estratégias posteriores receberam os casos que não atenderam à regra de parada, as diferenças foram interpretadas considerando essa seleção.

As propostas foram aplicadas em cópias de trabalho isoladas e avaliadas por reanálise estática e testes relacionados aos arquivos modificados. Na análise principal, o aceite exigiu a remoção do alerta original, a aprovação funcional na interpretação por critérios e a ausência de novos alertas, conforme as definições do Capítulo~\ref{cap:metodologia}. Essa interpretação distinguiu testes de comportamento legítimo daqueles que dependiam das vulnerabilidades intencionais da aplicação. A aplicação conjunta dos \textit{patches} selecionados também foi examinada, permitindo avaliar sua integração no projeto.

\textit{Prompts}, respostas, \textit{patches}, análises e testes foram registrados em arquivos estruturados. Esses dados sustentaram a comparação entre estratégias, modelos e tipos de alerta, incluindo custos calculados e tempos observados. As medições ausentes foram identificadas na análise, e o aceite foi interpretado como atendimento às verificações adotadas, sem constituir comprovação de segurança completa.

\section{Organização do Trabalho}
O restante deste trabalho está organizado da seguinte forma: o Capítulo 2 apresenta o referencial teórico e os trabalhos relacionados; o Capítulo~\ref{cap:metodologia} descreve a metodologia e os critérios de avaliação; o Capítulo~\ref{cap:implementacao} apresenta o desenvolvimento e a implementação do \textit{pipeline}; o Capítulo~\ref{cap:resultados} expõe os resultados, sua discussão e as respostas às perguntas de pesquisa; e o Capítulo~\ref{cap:consideracoes-finais} reúne as considerações finais e as propostas de trabalhos futuros.
